State James Theorem and its Lemmas


The lemmas leading to James Theorem are:
Let \(X\) be a Banach space and \(A\) a weakly closed nonempty subset of \(X.\) The following conditions are equivalent:
* \(A\) is weakly compact.
* For every \(f \in X^{\prime},\) there exists an element \(a_0 \in A\) such that \(\left|f\left(a_0\right)\right| = \sup_{a \in A} |f(a)|\).
* For any \(f \in X^{\prime}_{\mathbb{R}},\) there exists an element \(a_0 \in A\) such that \(f\left(a_0\right) = \sup_{a \in A} |f(a)|\).
* For any \(f \in X^{\prime}_{\mathbb{R}},\) there exists an element \(a_0 \in A\) such that \(f\left(a_0\right) = \sup_{a \in A} f(a)\).

A Banach space being reflexive if and only if its closed unit ball is weakly compact one deduces from this, 
since the norm of a continuous linear form is the upper bound of its modulus on this ball:

James's theorem
A Banach space \(X\) is reflexive if and only if for all \(f \in X^{\prime},\) 
there exists an element \(a \in X\) of norm \(\|a\| \leq 1\) such that \(f(a) = \|f\|\)


State the Theorem of Bishop \& Phelps


Let \(B \subseteq X\) be a bounded, closed, convex subset of a real [[Banach space]] \(X.\) Then the set of all continuous linear functionals \(f\) that achieve 
their supremum on \(B\) (meaning that there exists some \(b_0 \in B\) such that \(|f(b_0)| = \sup_{b \in B} |f(b)|\)) 
\[
    \left\{f \in X^\ast : f \text{ attains its supremum on } B\right\}
\]
is Dual norm-dense in the continuous dual space \(X^*\) of \(X.\) 

Importantly, this theorem fails for complex Banach spaces.
However, for the special case where \(B\) is the closed unit ball then this theorem does hold for complex Banach spaces.


Explain the concept of the weak derivative


For functions \( u, \phi \in H^1_1([a, b]), [a, b] \subset \mathbb{R} \) (Sobolev-Space) integration by parts states
\[
    \int_a^b u \phi' dx = \left [u \phi]_a^b - \int_a^b u' \phi dx \right.
\]
The weak derivative of a function \( u \in L^p([a, b]) \) (so only integrable but not cont. differentiable) is a function 
\( v \in L^p([a, b]) \) which mimics this rule:
\[
    \int_a^b u \phi' dx = \left [u \phi]_a^b - \int_a^b v' \phi dx \right.
\]
This is dependent on the choice of test functions \( \phi \) (and \( v \) should fit for all \( \phi \) in the choosen family), 
typically these are bump functions \( C^\infty_{c}([a, b]) \) and alike, hence \( \phi(a) = \phi(b) = 0 \), so the formula simplifies 
and we have in general for locally integrable function \( L^1_{\text{loc}}(U) \) and \( U \subseteq \mathbb{R}^n \)
and multi-index \( \alpha \):
\[
    \int_U u D^\alpha \phi = (-1)^{|\alpha|} \int_U v \phi.
\]

Example:
The characteristic function of the rational numbers \( \mathbf{1}_{\mathbb{Q}} \) with integration in respect to the Lebesgue measure
satisfies
\[
    \int \mathbf{1}_{Q} \phi d\lambda = 0
\]
hence \( v = 0 \) is a weak derivative of \( \mathbf{1} \).



Explain the weak formulation of PDE's and its relation to the Lax-Milgram Theorem

Let \(V\) be a [[Banach space]], let \(V'\) be the [[dual space]] of \(V\), let \(A\colon V \to V'\) be a [[linear map]], and let \(f \in V'\). 
A vector \(u \in V\) is a solution of the equation

\[ Au = f \]

if and only if for all \(v \in V\),

\[(Au)(v) = f(v).\]
A particular choice of \(v\) is called a test vector (in general) or a test function (if \(V\) is a function space).

To bring this into the generic form of a weak formulation, find \(u\in V\) such that

\[
    a(u,v) = f(v) \quad \forall v \in V,
\]

by defining the bilinear form (e.g. a matrix in the finite dim. case):

\[
    a(u,v) := (Au)(v).
\]

Let \(V\) be a real [[Hilbert space]] and \(a( \cdot ,\cdot )\) a [[bilinear form]] on {{nowrap|\(V\),}} which is
* bounded: \(|a(u,v)| \le C \|u\| \|v\|\) and
* coercive: \(a(u,u) \ge c \|u\|^2\)

Then, for any bounded \(f\in V'\), there is a unique solution \(u\in V\) to the equation
\[
    a(u,v) = f(v) \quad \forall v \in V
\]
and it holds
\[
    |u\| \le \frac1c \|f\|_{V'}
\]

To solve Poisson's equation
\(\nabla^2 u = f,\)

on a domain \(\Omega\subset \mathbb R^d\) with \(u=0\) on its boundary (topology), 
and to specify the solution space \(V\) later, one can use the \(L^2\) scalar product

\[
    \langle u,v\rangle = \int_\Omega uv\,dx
\]
to derive the weak formulation. Then, testing with differentiable function \(v\) yields

\[
    \int_\Omega ( \nabla^2 u ) v \,dx = \int_\Omega fv \,dx.
\]

The left side of this equation can be made more symmetric by integration by parts using Green's identity 
and assuming that \(v=0\) on \(\partial\Omega\):

\[
    \int_\Omega \nabla u \cdot \nabla v \,dx = \int_\Omega f v \,dx.
\]

This is what is usually called the weak formulation of Poisson's equation.

The relation to Lax-Milgram's Theorem is:
Here, choose \(V = H^1_0(\Omega)\) with the norm
\[
    \|v\|_V := \|\nabla v\|,
\]

where the norm on the right is the \(L^2\)-norm on \(\Omega\) (this provides a true norm on \(V\) by the Poincaré inequality.
But, we see that \(|a(u,u)| = \|\nabla u\|^2\) and by the Cauchy–Schwarz inequality, \(|a(u,v)| \le \|\nabla u\|\,\|\nabla v\|\).

Therefore, for any \(f \in [H^1_0(\Omega)]'\), there is a unique solution \(u\in V\) of Poisson's equation and we have the estimate
\[
\|\nabla u\| \le \|f\|_{[H^1_0(\Omega)]'}
\]


More generally for a PDE operator defined on some open \( U \subset \mathbb{R}^n \)
\[
P(x, \partial)u(x)=\sum a_{\alpha_1, \alpha_2, \dots, \alpha_n}(x) \, \partial^{\alpha_1}\partial^{\alpha_2}\cdots \partial^{\alpha_n} u(x), 
\]
The PDE
\[
P(x, \partial)u(x) = 0
\]
can be transformed into the weak form by multiplying and integrating with a test function
\[
\int_W u(x) Q(x, \partial) \varphi (x) dx =0,
\]
where 
\[
    Q(x, \partial)\varphi (x) = \sum (-1)^{| \alpha |} \partial^{\alpha_1} \partial^{\alpha_2} \cdots \partial^{\alpha_n} \left[a_{\alpha_1, \alpha_2, \dots, \alpha_n}(x)  \varphi(x) \right]
\]
Formally \( Q(x, \partial) \) is formal adjoint of \( P(x, \partial) \).


What is the space of distributions

The spaces of test functions and distributions are TVSs that are used in the definition and application of distributions. 
Test functions are usually \( C^\infty_c(U) \) for \( U \subset \mathbb{R}^n  \) (i.e. smooth functions with compact support in \( U \)). 
One is often interested in some restrictions e.g. \( C^k_c(U) \) and we endow the space with a certain topology, 
called the canonical LF-topology (locally convex inductive limit of a countable many Fréchet spaces), that makes 
\( C_{\text{c}}^{k}(U) \) into a complete Hausdorff locally convex TVS. The strong dual space of 
\( C_{\text{c}}^{k}(U) \) is called the \textbf{space of distributions} on  \( U \)
and is denoted by \( \mathcal{D}^{\prime }(U) \), endowed with the strong dual topology.

The LF-topology is quite technical and is the locally convex colimit of countable many Fréchet spaces
and it is a final topology since they contain each other.
For the distributions these are given by choosing a chain of compact sets \( \{K_i\}_{i \in \mathbb{N}} \) s.t. \( \cup_{i \in \mathbb{N}} K_i = \mathbb{R}^n \)
and taking the colimit of \( \{ C^\inf_c(K_i) \}_{i \in \mathbb{N}} \). It has a metric.
The strong dual topology is as usual the topology of uniform convergence on bounded subsets of \( C^\infty_c(U) \).

Sometimes one is interested more in well-behaved subspaces of \( \mathcal{D}^{\prime'}(U) \). 
One common example is the space of Schwarz functions/tempered distributions, which have Fourier Transformations. 

Proposition: If \( T \) is a linear functional on \( C_c^\infty(U) \) then the \( T \) is a distribution if and only if the following equivalent conditions 
for all compact \( K \subseteq U \) are satisfied
1. There exist constants 
\( |T(f)|\leq C\sup\{|\partial ^{\alpha }f(x)|:x\in U,|\alpha |\leq N\}.)
2. \( |T(f)|\leq C\sup\{|\partial ^{\alpha }f(x)|:x\in K,|\alpha |\leq N\}. \)
3. If for any sequence \( {\partial ^{\alpha }f_{i}\}_{i=1}^{\infty }} \) converges uniformly to zero on \( K \) for all multi-indices \( \alpha \), then 
\(T(f_{i})\to 0.\)

Further equivalent definitions are that 
1. \( T \) is cont.
2. \( T \) is uniformly cont.
3. \( T \) is a bounded operator
4. \( T \) is seq. cont.
A distribution is the further restriction 

For compact \( K \) and \( k = \infty \) \( C^k_c(K) \) has no metric. If \( k < \infty \) the space is Banach and has a norm/metric, given by 
\[
    |f| = \sup_{|p| < k} \left \sup_{x_0 \in K} | \partial^p  f(x_0) | \right)
\]
For open \( U \), \( C^k_c(U) \) has never a norm.

One can also define the transpose of a distribution, which allows to identify all Radon measures as distributions of \( C^0_c(U) \).
\( D^{\prime}(U) \) as the the dual of \( C^k_c(U) \) makes all linear partial differential operators of order \( k \) to distributions.


Give some typical geometric classifications of sets in TVS and its relation to NHDs bases

A subset \( A \) of a vector space is:
1. convex, if it includes the line segment joining any pair of its points.
2. absorbing (or radial) if for any \( x \) some multiple of \( A \) includes the line segment joining \( x \) and zero. 
That is, if there exists some \( \alpha_0 > 0 \) satisfying \( \alpha x \in A \) for every \( 0 \leq \alpha \leq \alpha_0 \)
Equivalently, \( A \) is absorbing if for each vector \( x \) there exists some \( \alpha_0 > 0 \) such that \( \alpha x\in A \) whenever \( - \alpha_0 \leq \alpha \leq \alpha_0 \)
3. circled (or balanced) if for each \( x \in A \) the line segment joining \( x \) and \( -x \) lies in A. 
That is, if for any \( x \in A \) and any \( |\alpha| \lq 1 \), we have \( \alpha x \in A \).
4. symmetric if \( x \in A \) implies \( - x \in A \).
5. star-shaped about zero if it includes the line segment joining each of its points with zero. 
That is, if for any \( x \in A \) and any \( 0 \leq \alpha \leq 1 \) we have \( \alpha x \in A \).

Structure Theorem If \( (X, \tau) \) is a tvs, then there is a neighborhood base \( \mathcal{B} \) at zero such that:
1. Each \( V \in \mathcal{B} \) is absorbing.
2. Each \( V \in \mathcal{B} \) is circled.
3. For each \( V \in \mathcal{B} \) there exists some \( W \in \mathcal{B} \) with \( W + W \subset V \).
4. Each \( V \in \mathcal{B} \) is closed.
Conversely, if a filter base \( \mathcal{B} \) on a vector space X satisfies properties \( (1-3) \), 
then there exists a unique linear topology \( \tau \) on X having \( \mathcal{B} \) as a neighborhood base at zero.

A Hausdorff topological vector space is metrizable if and only
if zero has a countable neighborhood base. In this case, the topology is generated
by a translation invariant metric.



Define dual pairs and dual topologies induced by subspaces.

A dual pair (or a dual system) is a pair \( \langle X, X' \rangle \) of vector
spaces together with a bilinear functional  \( (x, x') \mapsto \langle x, x' \rangle \) from \( X \times X' \to \mathbb{R} \), that
separates the points of \( X \) and \( X' \). That is:
1. The mapping \( \langle \cdot, \cdot, \rangle \) is bilinear.
2. Fix \( x \) in, if \( \langle x, x', \rangle = 0 \), for all \( x' \in X' \) then \( X = 0 \) and the symmetric version for \( x' \) fixed.

Since we can consider X to be a vector subspace of \( \mathbb{R}^X \), \( X \) inherits the product
topology of \( \mathbb{R}^X \) . This topology is referred to as the weak topology on \( X \) and is
denoted by \( \sigma(X, X') \), or simply \( w \). Since the product topology on \( \mathbb{R}^X \) is a locally
convex Hausdorff topology, the weak topology \( \sigma(X, X') \) is likewise Hausdorff and locally convex.
For this reason the weak topology is also known as the topology of pointwise convergence on X  . 

A family of seminorms that generates the weak
topology \( \sigma(X, X') \) is 
\[
  p_{x'}(x) = | \langle x, x' \rangle|, \quad x \in X
\]
The locally convex Hausdorff topology \( \sigma(X', X) \) is defined in a similar manner
and is the \( w^\ast \)-topology with the generating seminorms
\[
  p_{x}(x') = | \langle x, x' \rangle|, \quad x' \in X'
\]

Theorem 
If \( (X, X') \)is a dual pair, then the topological dual of the tvs \( X \), \( \sigma(X, X') \) is \( X' \). 
That is, if \( f : X \to \mathbb{R} \) is a \( \sigma(X, X') \)-continuous linear functional, then there exists a unique \( x' \in X' \) such that
\[
    f(x) = \langle x, x' \rangle
\]
